% -------------------------------------------------------
%  Section 12: Conclusion
% -------------------------------------------------------
\section{جمع‌بندی، محدودیت‌ها و کار‌های آینده}
\label{sec:conclusion}

\subsection{جمع‌بندی درباره مقاله}
مقاله \lr{Nie} ایده‌ای ساده و مفهومی تمیز پیشنهاد می‌کند: بخش‌های فایل اجرایی به کانال‌های جداگانه نگاشته شوند تا مرزبندی بخش‌ها به‌جای استنتاج از بافت، در ساختار ورودی کدگذاری شود. سنجه ارزیابی بیرونی، ثابت نگه داشتن عرض پایه در طرح‌های قطعه‌بندی و بخش محدودیت‌های صادقانه، امتیاز‌های واقعی این کار‌اند.

در مقابل، شواهد بار ادعا‌ها را تحمل نمی‌کنند: رتبه‌بندی ۲۰ پیکربندی بدون تکرار یا بازه اطمینان گزارش شده در حالی که پنج ردیف برتر در بازه‌ای به پهنای ده درصد نسبی جای می‌گیرند؛ دو معماری با چهار ابرپارامتر متفاوت مقایسه شده‌اند؛ و مهم‌تر از همه، ادعای محوری مقاومت در برابر مبهم‌سازی که انگیزه اصلی کل روش است در هیچ آزمایشی سنجیده نمی‌شود.

\subsection{جمع‌بندی درباره کار ما}
بازتولید ۲۰ پیکربندی روی \lr{BIG 2015} نشان داد طرح‌های خام و \lr{Mask} همراه تصویر خام با فاصله میانگین کمتر از \num{۰٫۰۰۷} زیان لگاریتمی بازتولید می‌شوند، ولی \lr{S4} و \lr{S5} بدون تصویر خام با فاصله‌ای یک مرتبه بزرگی بیشتر واگرا می‌شوند. ترتیب طرح‌های عرض روی \lr{ResNet50} دقیقاً بازتولید شد ولی روی \lr{VGG16} برعکس درآمد، که شکنندگی این رتبه‌بندی‌ها را تأیید می‌کند.

انتقال روش به \num{۲۷٬۵۷۱} سفت‌افزار \lr{ARM/Zephyr} پاسخ روشنی به پرسش انتقال‌پذیری داد. مقاله هرگز ادعا نکرده که قطعه‌بندی بهترین صحت را می‌دهد (در جدول ۵ خود مقاله هم \lr{S3} بهترین است)؛ ادعا این است که قطعه‌بندی به صحتی \emph{نزدیک} می‌رسد و در عوض استقلال از چیدمان فضایی را به دست می‌آورد. در دامنه نهفته هزینه این معامله بسیار سنگین‌تر می‌شود: بهترین \lr{S4} به \pct{۲۹٫۰۳} و بهترین \lr{S5} بدون تصویر خام به \pct{۲۸٫۵۸} می‌رسند، در حالی که یک رستر ساده با عرض ثابت \pct{۶۵٫۷۵} می‌دهد. تنها خانواده‌ای که نزدیک می‌ماند همان است که تصویر خام را نگه می‌دارد و در نتیجه همان وابستگی فضایی را بازمی‌گرداند که قرار بود حذف شود. در مقابل، سهم دوم مقاله یعنی اهمیت هم‌ترازی عرض، در این دامنه به‌مراتب قوی‌تر از دامنه اصلی تأیید شد: فاصله عرض ثابت و عرض تطبیقی از \num{۰٫۰۰۲} زیان لگاریتمی روی \lr{BIG 2015} به ۳۵ واحد صحت روی \lr{ARM} می‌رسد. مشاهده مقاله درباره سهم نابرابر بخش‌ها نیز تأیید شد، هرچند با فهرست کاملاً متفاوتی از بخش‌های مؤثر.

بازبینی انتقادی کار خودمان دو یافته داشت. نخست، پروتکل ارزیابی اولیه سنجه‌ها را روی کل مجموعه‌داده محاسبه می‌کرد و سرخط \pct{۷۹٫۸۷} در واقع ترکیب $0.8 \times 83.52 + 0.2 \times 65.26$ بود؛ پس از اصلاح، هیچ‌یک از دوازده پیکربندی مشترک با چرخه پایه از نظر صحت کلی بهبود نشان نمی‌دهد. دستاورد واقعی چرخه سوم جای دیگری است: کلاس \lr{benign} که در چرخه پایه با بازخوانی صفر فرو ریخته بود، با نمونه‌برداری وزن‌دار به بازخوانی \pct{۶۰٫۳} بازگشت و \lr{F1} کلان بهترین مدل از \num{۰٫۵۶۸۴} به \num{۰٫۶۰۸۴} رسید؛ برای یک سامانه تشخیص این تفاوت میان مدل بی‌استفاده و مدل قابل بحث است. دوم، تفکیک تصادفی در سطح نمونه نشت گروهی دارد و حتی عدد اصلاح‌شده نیز سقفی خوش‌بینانه است.

\subsection{کار‌های آینده}
سه اولویت، به ترتیب اهمیت:

\begin{enumerate}[leftmargin=1.6em,itemsep=0.15em,topsep=0.25em]
    \item \textbf{ارزیابی با تفکیک گروهی.} بازآموزی پیکربندی‌های برتر با گروه‌بندی بر پایه شناسه ساخت پایه و سپس شناسه دستور تزریق؛ تنها راه سنجش تعمیم واقعی به سفت‌افزار و گونه دیده‌نشده.
    \item \textbf{آزمون مقاومت در برابر مبهم‌سازی.} ارزیابی مدل‌های آموزش‌دیده روی نمونه‌هایی با بخش‌های بازآراییده یا داده زائد تزریق‌شده و مقایسه افت \lr{S3} با \lr{S4}؛ تنها راه سنجش مستقیم ادعای محوری مقاله.
    \item \textbf{اصلاح سر ادغام هرمی.} پیاده‌سازی درست \lr{PSPNet} با الحاق مکانی نقشه‌های ادغام‌شده و نرخ یادگیری جداگانه برای سر تازه؛ تنها پس از آن می‌توان درباره سودمندی ادغام هرمی داوری کرد.
\end{enumerate}

در کنار این سه، تکرار اجرا‌ها با چند بذر کمترین کار لازم برای قابل دفاع کردن هر رتبه‌بندی است.
